% Vietnamese translation for OI Public Library
% Source-File: IOI中国国家候选队论文/2009/漆子超/分治算法在树的路径问题中的应用.ppt
\documentclass[11pt]{article}
\usepackage{oipl-vn}

\newcommand{\Oh}{\mathcal{O}}

\title{Ứng dụng thuật toán chia để trị vào bài toán đường đi trên cây\\{\large Ghi chú theo slide}}
\author{Qi Zichao\\Trường Trung học Yali, Trường Sa}
\date{Luận văn đội tuyển dự tuyển quốc gia IOI 2009}

\begin{document}
\maketitle

\origtitle{Divide and conquer for path problems on trees}
\sourcepath{IOI-China-Candidate-Team-Papers/2009/Qi-Zichao/tree-path-divide-conquer-slides.ppt}

\section{Mạch trình bày}

Slide so sánh hai cách xử lý đường đi trên cây: chia để trị trên cây và phân
rã đường. Cả hai đều khai thác việc một đường đi bất kỳ có thể được phân tách
qua một điểm, một cạnh hoặc một số đoạn có cấu trúc.

\section{Chia để trị trên cây}

Trong chia để trị theo đỉnh, ta chọn một trọng tâm, xử lý mọi đường đi đi qua
trọng tâm, rồi đệ quy trên các thành phần còn lại. Nếu chọn trọng tâm tốt,
độ sâu đệ quy là \(\Oh(\log n)\). Với chia để trị theo cạnh, ý tưởng tương tự
nhưng điểm tách là cạnh, phù hợp với một số công thức gộp khác.

\section{Phân rã đường}

Phân rã nặng nhẹ chia đường đi giữa hai đỉnh thành một số đoạn trên các chuỗi
nặng. Nhờ vậy truy vấn và cập nhật đường đi có thể chuyển thành các thao tác
trên mảng hoặc cây đoạn.

\section{Các ví dụ}

\textit{Đếm cặp đỉnh trên cây}, \textit{Free Tour 2}, \textit{Query On a
Tree}, \textit{Cây đen trắng} và \textit{Query On a Tree IV} minh họa cách
chọn kỹ thuật theo dạng truy vấn: đếm toàn cục thường hợp với chia để trị,
còn truy vấn động nhiều lần thường hợp với phân rã đường.

\end{document}
